\documentclass[11pt,a4paper]{article}
\usepackage[T1]{fontenc}
\usepackage[utf8]{inputenc}
\usepackage{lmodern,microtype}
\usepackage[margin=26mm]{geometry}
\usepackage{amsmath,amssymb,amsthm,mathtools}
\usepackage{booktabs,array,longtable,enumitem,xcolor}
\usepackage[colorlinks=true,linkcolor=blue!45!black,citecolor=blue!45!black,urlcolor=blue!45!black]{hyperref}
\usepackage{fancyhdr}
\pagestyle{fancy}\fancyhf{}
\fancyhead[L]{\small Stern root certificates and a bilinear Hankel identity}
\fancyhead[R]{\small Research round T86--T90}
\fancyfoot[C]{\thepage}\setlength{\headheight}{14pt}
\setlength{\parskip}{4pt}\setlength{\parindent}{1.1em}
\setlist{itemsep=3pt,topsep=5pt}
\newtheorem{theorem}{Theorem}[section]
\newtheorem{lemma}[theorem]{Lemma}
\newtheorem{proposition}[theorem]{Proposition}
\newtheorem{corollary}[theorem]{Corollary}
\theoremstyle{definition}\newtheorem{definition}[theorem]{Definition}
\theoremstyle{remark}\newtheorem{remark}[theorem]{Remark}
\newcommand{\Z}{\mathbb Z}\newcommand{\Q}{\mathbb Q}
\newcommand{\C}{\mathbb C}\newcommand{\R}{\mathbb R}
\hypersetup{pdftitle={Finite Certificates for Stern Roots and a Bilinear Hankel Identity: Research Round T86-T90},pdfauthor={KEIO}}
\title{Finite Certificates for Stern Roots\\and a Bilinear Hankel Identity\\[5pt]\large Complete arguments, auxiliary classifications, and a five-target audit}
\author{KEIO}
\date{26 October 2026\\[3pt]\small Research round T86--T90}
\begin{document}
\maketitle
\begin{abstract}
This manuscript records a five-target research round, separating exact
prior-question applications from auxiliary results and unsuccessful
attempts. A parity-separated quadratic envelope, combined with finite
polynomial positivity certificates, proves the entire three-family
assertion in Ulas's Conjecture~4.3(1): the known zero $-1$ is the
unique real zero and is simple. For Cigler's Conjecture~4.1,
we reduce boundary-deformed path Hankel determinants to a reciprocal
Cauchy determinant. Its spectral root equation identifies three classical
Miwa directions with the integer shifts $1,-k,k+1$, yielding the original
bilinear recurrence for every parameter and every permitted determinant
size. An exact classification of Laplacian peak transfer on uniform
two-level trees is included as an auxiliary result; the underlying
existence question has a prior public solution. The least-dimensional
Ehrhart counterexample question and the all-odd-period PT existence
conjecture remain unresolved here. No historically unprecedented
foundational principle, exhaustive priority determination, or five
independent open-problem solutions are claimed.
\end{abstract}

\section{The five targets and what is actually established}

The labels T86--T90 identify research tracks, not a promise of five
successful theorems. We require an unchanged, previously stated target
before counting an application as a prior-question resolution. Supporting
lemmas, already known solutions, special cases, and unsuccessful
constructions are not added to that count.

\begin{center}\small
\begin{tabular}{@{}>{\raggedright\arraybackslash}p{.08\textwidth}>{\raggedright\arraybackslash}p{.39\textwidth}>{\raggedright\arraybackslash}p{.45\textwidth}@{}}
\toprule
Track & Original target & Outcome in this manuscript \\
\midrule
T86 & Clarke Question~6.8: least dimension of an ineffective polynomial equivariant numerator. & Unresolved. A concrete candidate is effective; it is not a counterexample. \\
T87 & Faust et al.\ Conjecture~2: nonzero PT-symmetric zero-isospectral potentials for every period greater than three. & Unresolved. A complete obstruction eliminates one sparse construction for all odd periods at least seven. \\
T88 & Ulas Conjecture~4.3(1): exactly one real root in each of three original Stern families. & Complete proof for all three families and all $n\geq2$, including simplicity and positivity of the quotient. \\
T89 & Coutinho--Guo--Schmeits Open Problem~7.1: infinitely many non-star trees with Laplacian peak transfer. & Excluded from new-resolution counting because of prior solution. A full classification inside a two-parameter tree family is proved. \\
T90 & Cigler Conjecture~4.1: an all-column bilinear Hankel recurrence. & Complete proof of the original recurrence, including exceptional parameters and vanishing determinants. \\
\bottomrule
\end{tabular}
\end{center}

The principal mathematical claims are the unchanged Ulas clause and
Cigler recurrence. Neither result is a proof of every conjecture in its
source. The methods use classical Stern identities, concavity, Bernstein
positivity, Christoffel determinants, and the Hirota--Miwa identity.
The additional content proposed here consists of the finite certificates
for the specific infinite Stern families and the exact reciprocal-root
reduction of the Hankel problem. These applications do not establish a
historically new foundational principle. A bounded literature search
cannot prove that no earlier identical solution exists.

All statements presented as theorems below have complete unconditional
arguments. This is an internally checked research draft, not an externally
refereed paper or a proof-assistant development. The user's full set of
historical success conditions is therefore not certified.


\section{Finite envelopes and the real-root conjecture for three Stern families}
\label{sec:stern-three-families}

Let $B_0(t)=0$, $B_1(t)=1$, $B_{2m}(t)=tB_m(t)$, and
$B_{2m+1}(t)=B_m(t)+B_{m+1}(t)$. The original target is
Conjecture~4.3(1) of Ulas~\cite{Ulas20}, for the unchanged indices
\begin{align*}
 s_{0,n}&=2^{2n+4}-9\cdot2^{n+1}-1,\\
 s_{2,n}&=2^{2n+8}-9\cdot2^{n+4}-13,\\
 s_{3,n}&=2^{2n+8}-51\cdot2^{n+2}-1.
\end{align*}
That numbered assertion says that each $B_{s_{i,n}}$ has exactly one real
root for $i\in\{0,2,3\}$ and every $n\geq2$. We prove the stronger
strict positivity of the quotient by its known factor $t+1$.
This is different from the eventual-monotonicity counterexample in T81.
It does not assert the irreducibility asked for in Conjecture~4.3(2),
or settle the other clauses of Conjecture~4.3.

\subsection{The auxiliary finite-certificate problem}

The auxiliary problem is to certify the sign of an infinite family
\[
 F_n(x)=a(x)+b(x)(-x)^n+c(x)x^{2n},\qquad n\geq2,
\]
without bounding $n$. Splitting by parity turns the moving power into
one real variable in a known interval. The following observation gives a
finite reduction when the resulting quadratic is concave.

\begin{lemma}[Parity-separated quadratic envelopes]\label{lem:stern-envelope}
Fix $x>0$ and put
$g_e(v)=a(x)+(-1)^e b(x)v+c(x)v^2$ for $e\in\{2,3\}$.
For $n\geq2$ of the same parity as $e$, we have $F_n(x)=g_e(x^n)$.
If $0<x<1$ and $c(x)\leq0$, then
\[
 F_n(x)\geq\min\{a(x),F_e(x)\}.
\]
If $x>1$, $c(x)<0$, and $g'_e(x^e)<0$, then
\[
 F_n(x)\leq F_e(x).
\]
\end{lemma}
\begin{proof}
In the first case, $0<x^n\leq x^e$ and a concave quadratic on
$[0,x^e]$ takes its minimum at an endpoint. In the second case,
$x^n\geq x^e$ and $g'_e(v)=(-1)^e b(x)+2c(x)v$ is strictly decreasing
and already negative at $x^e$.
\end{proof}

The remaining finite positivity assertions below use the elementary
Bernstein basis
\[
 \mathcal B_{j,N}(u)=\binom Nj u^j(1-u)^{N-j},\qquad 0\leq j\leq N.
\]
If $p(ru)=D^{-1}\sum_{j=0}^N v_j\mathcal B_{j,N}(u)$ with
$D>0$, $v_j\geq0$, and $v_0,v_N>0$, then $p(x)>0$ on $[0,r]$.
Indeed all basis functions are positive on $(0,1)$, and the endpoint
values are $v_0/D$ and $v_N/D$. This proof uses only finite polynomial
identities; no numerical root isolation is assumed.

\subsection{Exact family formulas and finite positivity data}

Put $A(x)=2x^3+2x^2-3x+1$ and write
\[
 F_{i,n}(x):=(1+x)^2B_{s_{i,n}}(-x)
        =a_i(x)+b_i(x)(-x)^n+c_i(x)x^{2n}.
 \tag{S1}\label{eq:stern-S1}
\]
The coefficient polynomials are
\begin{align*}
 a_0&=A,\qquad b_0=-2(1-x)^2,\qquad c_0=-x(x^2-x+2),\\
 a_2&=(2x^2-2x+1)A,\qquad c_2=x^2(2-3x)(x^2-x+2),\\
 b_2&=2x(1-x)^2(x^3+x^2-3x+2),\\
 a_3&=2x^5-4x^3+6x^2-3x+1,\\
 b_3&=2x(1-x)(3x^2-2x+1),\qquad c_3=-x^3(5x^2-5x+2).
\end{align*}
Here and below the argument $x$ on the right is suppressed when
unambiguous. For completeness, these identities follow from the Stern
addition formula
\[
 B_{m2^a-r}=B_{2^a-r}B_m+B_rB_{m-1}
 \qquad(0\leq r\leq2^a),
\]
and its plus version. A convenient way to evaluate the complementary
indices is, for $2^b\geq r$ and $a\geq b$,
\[
 B_{2^a-r}(t)=
 \frac{t^{a-b}\{B_r(t)+(t-1)B_{2^b-r}(t)\}-B_r(t)}{t-1}.
 \tag{S2}\label{eq:stern-S2}
\]
The plus formula with
$2^a-r=2^b(2^{a-b}-1)+(2^b-r)$ proves \eqref{eq:stern-S2} directly;
all displayed quotients are polynomials.
For $i=0,2,3$, respectively, apply the minus formula to
\[
 s_{0,n}=2^{n+1}(2^{n+3}-9)-1,\quad
 s_{2,n}=2^{n+4}(2^{n+4}-9)-13,\quad
 s_{3,n}=2^{n+2}(2^{n+6}-51)-1.
\]
Then use \eqref{eq:stern-S2} with $r=9,10,13,51,52$ and
$b=4,4,4,6,6$. The small polynomials needed are
\[
 B_9=(t+1)^2,\ B_{10}=t+2t^2,\ B_{13}=1+2t+2t^2,
 \quad B_{51}=1+3t+5t^2+3t^3,\quad B_{52}=t^2+2t^3+2t^4.
\]
Substitution and multiplication give \eqref{eq:stern-S1} for every
$n\geq2$, rather than merely a polynomial interpolation in $n$.

For $e=2,3$ define the following six base quotients by the identities
\[
 F_{i,e}(x)=(1-x)(1+x)^2q_{i,e}(x).
 \tag{S3}\label{eq:stern-S3}
\]
Explicitly,
\begin{align*}
q_{0,2}&=x^4-2x^3+5x^2-4x+1,\\
q_{0,3}&=x^6-2x^5+5x^4-6x^3+7x^2-4x+1,\\
q_{2,2}&=3x^6-10x^5+23x^4-26x^3+17x^2-6x+1,\\
q_{2,3}&=3x^8-8x^7+21x^6-32x^5+37x^4-30x^3+17x^2-6x+1,\\
q_{3,2}&=5x^6-10x^5+17x^4-16x^3+11x^2-4x+1,\\
q_{3,3}&=5x^8-10x^7+17x^6-22x^5+23x^4-18x^3+11x^2-4x+1.
\end{align*}
All six are strictly positive on $[0,\infty)$. To give finite exact
certificates, first note
\[
 q_{0,2}(x)=(x^2-x)^2+(2x-1)^2>0\qquad(x\in\mathbb R).
\]
For the other five, the following table records $(N,D,(v_j)_{j=0}^N)$
in the identity
\[
q(x)=D^{-1}\sum_jv_j\mathcal B_{j,N}(x).
\]
\begin{center}\small
\begin{tabular}{@{}lrrl@{}}
\toprule
Polynomial & $N$ & $D$ & $(v_0,\ldots,v_N)$ \\
\midrule
$q_{0,3}$ & 6 & 30 & $(30,10,4,3,8,20,60)$\\
$q_{2,2}$ & 6 & 30 & $(30,0,4,3,4,10,60)$\\
$q_{2,3}$ & 8 & 140 & $(140,35,15,5,4,5,20,70,420)$\\
$q_{3,2}$ & 6 & 15 & $(15,5,6,6,10,15,60)$\\
$q_{3,3}$ & 8 & 140 & $(140,70,55,50,56,65,90,105,420)$\\
\bottomrule
\end{tabular}
\end{center}
These identities prove positivity on $[0,1]$. On $[1,\infty)$, the
coefficient vectors of $q(1+y)$, in ascending powers of $y$, are
\begin{center}\small
\begin{tabular}{@{}ll@{}}
\toprule
Polynomial & Coefficients of $q(1+y)$ \\
\midrule
$q_{0,2}$ & $(1,4,5,2,1)$\\
$q_{0,3}$ & $(2,8,14,14,10,4,1)$\\
$q_{2,2}$ & $(2,10,22,26,18,8,3)$\\
$q_{2,3}$ & $(3,20,60,106,122,94,49,16,3)$\\
$q_{3,2}$ & $(4,18,40,52,42,20,5)$\\
$q_{3,3}$ & $(3,18,60,124,168,150,87,30,5)$\\
\bottomrule
\end{tabular}
\end{center}
Their positive coefficients prove the remaining assertion.

We also need $a_i>0$ on $[0,1]$. For $a_0=A$ and $a_3$, their
Bernstein certificates have degree five, denominator five, and vectors
\[
 A:\ (5,2,0,0,3,10),\qquad
 a_3:\ (5,2,2,3,3,10).
\]
The factor $2x^2-2x+1=2(x-1/2)^2+1/2$ proves the assertion for $a_2$.
Finally put
\[
 J(x)=a_2(x)-b_2(x)x^3.
\tag{S4}\label{eq:stern-S4}
\]
The exact degree-nine Bernstein certificate for $J(2u/3)$ has
denominator $413343$ and vector
\[
\begin{split}
(&413343,260253,158193,95499,57915,\\
 &36369,26973,28623,47159,97125).
\end{split}
\]
Thus $J(x)>0$ for $0\leq x\leq2/3$. Every entry in these tables is
rational or integral; their defining identities can be verified by
expansion and form part of the proof.

\subsection{The full original real-root assertion}

\begin{theorem}[Ulas Conjecture 4.3(1), with a positive quotient]
\label{thm:ulas-three-root-families}
For every $n\geq2$ and $i\in\{0,2,3\}$,
\[
 \frac{B_{s_{i,n}}(t)}{t+1}>0\qquad(t\in\mathbb R),
\]
where the quotient is the polynomial obtained by division by $t+1$.
Consequently $-1$ is the unique real zero of $B_{s_{i,n}}$, and that
zero is simple. This proves the whole assertion of
\cite[Conjecture~4.3(1)]{Ulas20}.
\end{theorem}
\begin{proof}
For $t\geq0$, the Stern recurrence gives nonnegative coefficients and
constant coefficient one at these odd indices, so the asserted quotient
is positive. For $t=-x<0$, it suffices to prove
$F_{i,n}(x)/(1-x)>0$, treating $x=1$ by its removable value.

Suppose first $0<x<1$. For $i=0,3$ we have $c_i(x)<0$. For each
parity take $e=2$ or $e=3$ in Lemma~\ref{lem:stern-envelope}.
Both endpoints are positive: $a_i(x)>0$ by the certificates, and
$F_{i,e}(x)=(1-x)(1+x)^2q_{i,e}(x)>0$. Thus $F_{i,n}(x)>0$.
For $i=2$ the same argument applies on $2/3\leq x<1$, because
$c_2\leq0$ there. On $0<x\leq2/3$, we have $c_2\geq0$ and
$b_2\geq0$, since $x^3+x^2-3x+2\geq 2-3x\geq0$.
If $n$ is even, all three terms in \eqref{eq:stern-S1} are nonnegative
and $a_2>0$. If $n$ is odd, then $n\geq3$, $x^n\leq x^3$, and
\[
 F_{2,n}(x)=a_2-b_2x^n+c_2x^{2n}\geq a_2-b_2x^3=J(x)>0.
\]
This proves the required sign for every family on $(0,1)$.

Now suppose $x>1$. All three $c_i$ are negative. We check the
second condition in Lemma~\ref{lem:stern-envelope} at $e=2,3$.
For $i=0$, $b_0<0$, so the even derivative is negative. The odd one is
\[
 -b_0+2c_0x^3=2(x-1)^2-2x^4(x^2-x+2)<0.
\]
Here $x^2-x+2\geq2$ and $x^4>(x-1)^2$.
For $i=2$, $b_2>0$ on $x>1$, since
$x^3+x^2-3x+2=x(x^2-1)+(x-1)^2+1>0$. Thus the odd derivative is negative.
The even derivative is $-2xH_2(x)$, where
\[
 H_2=x^3(3x-2)(x^2-x+2)-(x-1)^2(x^3+x^2-3x+2).
\]
For $i=3$, $b_3<0$, so the even derivative is negative; the odd
one is $-2xH_3(x)$, where
\[
 H_3=x^5(5x^2-5x+2)-(x-1)(3x^2-2x+1).
\]
The positivity needed here follows from the exact expansions
\begin{align*}
 H_2(1+y)&=2+13y+30y^2+36y^3+24y^4+12y^5+3y^6,\\
 H_3(1+y)&=2+13y+46y^2+92y^3+110y^4+77y^5+30y^6+5y^7.
\end{align*}
Thus for the matching parity,
\[
 F_{i,n}(x)\leq F_{i,e}(x)
      =(1-x)(1+x)^2q_{i,e}(x)<0.
\]
This proves positivity of the quotient also when $x>1$.

It remains to verify division and strict positivity at $x=1$.
Formula~\eqref{eq:stern-S1} gives $F_{i,n}(1)=0$, proving divisibility
by $t+1$. Differentiation of that same identity gives the removable
values
\[
 \left.\frac{B_{s_{0,n}}(t)}{t+1}\right|_{t=-1}=n-1,\qquad
 \left.\frac{B_{s_{2,n}}(t)}{t+1}\right|_{t=-1}=n,\qquad
 \left.\frac{B_{s_{3,n}}(t)}{t+1}\right|_{t=-1}=n+1+(-1)^n.
\]
They are strictly positive for $n\geq2$. The quotient therefore never
vanishes on $\mathbb R$, and the root $-1$ is simple.
\end{proof}

\paragraph{Attribution and scope.}
The underlying Stern addition identities are prior. Concavity,
parity splitting, and Bernstein positivity are classical tools.
The proposed additional content is their finite-envelope application to
all three original families, including the convex subinterval of the
second family and the simple-root boundary values. No historically
unprecedented general principle is inferred from this application.
A bounded source search did not locate an earlier proof of this precise
three-family statement; this is not a guarantee of exhaustive priority.

\subsection{The fourth family: an exact classification and the multiplicity issue}
\label{sec:stern-fourth-family}

For completeness consider the remaining original family
\[
 s_{1,n}=2^{2n+5}-9\cdot2^{n+2}-5,\qquad n\geq1.
\]
Conjecture~4.3(3) of Ulas~\cite{Ulas20} says that
$B_{s_{1,n}}$ has exactly three real roots for $n\geq3$.
The exact classification below distinguishes distinct roots from roots
counted with multiplicity. In particular, $n=3$ has only one distinct
real root, of multiplicity three. We retain that endpoint rather than
quietly replacing the original range by $n\geq4$.

\begin{lemma}[A uniform factorization for the fourth family]
\label{lem:stern-fourth-factorization}
Put
\[
 M_n(t)=(t+2)t^n-2t-1,\qquad
 N_n(t)=(t^2+t+2)t^n+2t^3-2t^2-3t-1.
\]
For every $n\geq1$,
\[
 B_{s_{1,n}}(t)=\frac{M_n(t)N_n(t)}{(t-1)^2}.
 \tag{S5}\label{eq:stern-S5}
\]
Both quotients $M_n/(t-1)$ and $N_n/(t-1)$ are polynomials.
\end{lemma}
\begin{proof}
Write $s_{1,n}=2^{n+2}(2^{n+3}-9)-5$. The addition identity gives
\[
 B_{s_{1,n}}=B_{2^{n+2}-5}B_{2^{n+3}-9}
                 +B_5B_{2^{n+3}-10}.
\]
The complement identity \eqref{eq:stern-S2} yields
\[
 B_{2^{n+2}-5}=\frac{M_n}{t-1},\qquad
 B_{2^{n+3}-10}=\frac{tM_n}{t-1},\qquad
 B_{2^{n+3}-9}=\frac{t^n(t^2+t+2)-(t+1)^2}{t-1}.
\]
Since $B_5=1+2t$, their substitution proves \eqref{eq:stern-S5}.
This is a recurrence derivation, independent of the displayed
parity-dependent factors in the source.
\end{proof}

\begin{theorem}[Complete real-root classification of the fourth family]
\label{thm:stern-fourth-classification}
For $n\geq4$, $B_{s_{1,n}}$ has exactly three real roots, all simple:
\[
 -r_n^{-1},\qquad -1,\qquad -r_n,
 \qquad 0<r_n<1.
\]
If $n$ is even, then $0<r_n<1/2$; if $n$ is odd, then
$1/2<r_n<1$. The exceptional initial cases are
\begin{align*}
 B_{s_{1,1}}&=(t+1)(3t^2+2t+1),\\
 B_{s_{1,2}}&=(t+1)(t^2+3t+1)^2,\\
 B_{s_{1,3}}&=(t+1)^3(t^4+2t^3+6t^2+4t+1).
\end{align*}
Thus the numbers of distinct real roots for $n=1,2,3$ are respectively
$1,3,1$, with multiplicities respectively $(1)$, $(2,1,2)$, and $(3)$.
For every $n\geq3$ the number of real roots counted with multiplicity is
exactly three.
\end{theorem}
\begin{proof}
The three initial factorizations follow from
Lemma~\ref{lem:stern-fourth-factorization}. Their quadratic factor at
$n=1$ has negative discriminant. The quartic factor at $n=3$ is positive
on the real line: at $t=-x$ it is
\[
 x^4-2x^3+6x^2-4x+1
       =(x^2-x)^2+(2x-1)^2+x^2>0.
\]
The strictness follows because these three squares cannot vanish
simultaneously. The quadratic at $n=2$ has two distinct negative roots,
neither equal to $-1$.

Suppose first that $n\geq4$ is even. Then
\[
 B_{s_{1,n}}=(t+1)\,\frac{M_n}{t-1}\,\frac{N_n}{t^2-1}.
\]
We show that the last quotient is positive on the real line. For
$t=-x$ with $x>0$ and $x\ne1$, it equals
\[
 G_n(-x)=\frac{(x^2-x+2)x^n-A(x)}{x^2-1},
 \qquad A(x)=2x^3+2x^2-3x+1.
\]
Since $n\geq4$, comparison with $x^4$ shows $G_n(-x)\geq G_4(-x)$
on both $0<x<1$ and $x>1$. The base polynomial satisfies
\[
\begin{split}
 G_4(-x)&=x^4-x^3+3x^2-3x+1\\
 &=\left(x^2-\frac x2\right)^2
    +\frac{11}{4}\left(x-\frac6{11}\right)^2+\frac2{11}>0.
\end{split}
\]
At $x=1$, differentiation of numerator and denominator gives
$G_n(-1)=n-3>0$. For $t\geq0$, positivity also follows from
\[
 \frac{N_n(t)}{t-1}=B_{2^{n+3}-9}(t)+tB_5(t)>0,
\]
and $t+1>0$.

It remains to count zeros of $M_n/(t-1)$. This is the reciprocal
polynomial
\[
 \frac{M_n(t)}{t-1}=1+3(t+t^2+\cdots+t^{n-1})+t^n.
\]
At $t=-x$ it equals
\[
 \frac{1-2x-(2-x)x^n}{x+1}.
\]
For $0<x<1$, its numerator has derivative
\[
 -2-x^{n-1}\{2n-(n+1)x\}<0.
\]
It takes values $1$ at $x=0$, a negative value at $x=1/2$, and
$-2$ at $x=1$. Thus it has exactly one zero $r_n\in(0,1/2)$,
and that zero is simple. Reciprocity gives exactly one further zero
at $1/r_n>1$, also simple; $x=1$ is not a zero. There are no
nonnegative roots because the coefficients are positive. This proves
the even case, including simplicity of the factor $t+1$.

Now suppose $n\geq5$ is odd. Then
\[
 B_{s_{1,n}}=(t+1)\,\frac{M_n}{t^2-1}\,\frac{N_n}{t-1}.
\]
The last quotient is positive. Indeed, for $t=-x$ with $x>0$ it is
\[
 \frac{A(x)+(x^2-x+2)x^n}{x+1}>0.
\]
We already proved $A(x)>0$ on $[0,1]$; for $x=1+y\geq1$, the
identity $A(1+y)=2+7y+8y^2+2y^3$ proves the rest. At $t\geq0$ the
same Stern-polynomial expression as above proves positivity.

The remaining factor $K_n=M_n/(t^2-1)$ is reciprocal of even degree
$n-1$. This follows either by division or from
$t^{n+1}M_n(1/t)=-M_n(t)$ and the same anti-reciprocity of $t^2-1$.
At $t=-x$ it is
\[
 K_n(-x)=\frac{f_n(x)}{x^2-1},\qquad
 f_n(x)=2x-1-(2-x)x^n.
\]
On $0<x<1$,
\[
 f_n''(x)=-n x^{n-2}\{2(n-1)-(n+1)x\}<0.
\]
Furthermore $f_n(0)=-1$, $f_n(1)=0$,
$f_n'(0)=2$, and $f_n'(1)=3-n<0$. The strictly decreasing
first derivative therefore has exactly one zero. The function increases
to its unique maximum and then decreases to zero at $1$; since the
maximum is positive, there is exactly one interior zero, on the strictly
increasing part. It is simple. As $f_n(1/2)<0$, that zero $r_n$ lies
in $(1/2,1)$. Reciprocity supplies exactly one further simple zero
at $1/r_n$. At the removable point,
\[
 K_n(-1)=\frac{3-n}{2}<0,
\]
so $-1$ is not a zero of $K_n$. Finally, $K_n(t)>0$ for $t\geq0$:
$M_n/(t-1)$ has positive coefficients, and division by $t+1>0$ preserves
its positive values. This proves the odd case and completes the full
classification.
\end{proof}

\paragraph{Exact relationship with the original statement.}
Theorem~\ref{thm:stern-fourth-classification} proves the stated
three-root assertion on its full range $n\geq3$ when roots are counted
with multiplicity. Under a distinct-root interpretation, the endpoint
$n=3$ refutes it, and the theorem gives the exact corrected classification.
The source does not explicitly specify multiplicity in Conjecture~4.3(3),
although Conjecture~4.2 separately defines a count using set cardinality.
For that reason this is not counted as another unqualified resolution of
an unambiguous original conjecture.

The inspected arXiv v1, page~16, also has a coefficient discrepancy in its
odd-index factorization: in its first factor, the coefficient of
$t(t^{2(m+1)}-1)/(t-1)$ must be $4$, not $1$. For example $m=1$
already distinguishes the two formulas. Identity~\eqref{eq:stern-S5}
provides the corrected all-$n$ expression directly from the recurrence;
no conclusion depends on the printed factorization. This observation is
about the inspected preprint, not an uninspected journal revision.

\section{A boundary spectral constraint and Cigler's bilinear Hankel recurrence}
\label{sec:cigler-boundary}

The target is Conjecture~4.1 of Cigler's
\emph{Some remarks and conjectures about Hankel determinants}
\cite{CiglerRemarks}. Cigler also posed the identical recurrence in his
MathOverflow question of 18 October 2023 \cite{CiglerMO23}.
Let $a_{n,j}=a_{n,j}(x,c)$ be defined by
\begin{gather*}
 a_{0,j}=\delta_{0j},\qquad a_{n,-1}=0,\qquad
 s_0=c+x,\quad s_j=c\ (j>0),\\
 a_{n,j}=a_{n-1,j-1}+s_j a_{n-1,j}+a_{n-1,j+1}.
\end{gather*}
Write
\[
 D_{k,n}(x,c)=\det(a_{i+j,k}(x,c))_{0\leq i,j<n},
 \qquad D_{k,0}=1.
\]
The original conjecture is a three-term bilinear identity for every
$k\geq1$ and $n\geq2k+2$. The auxiliary problem developed below asks
which spectral constraints make a Cauchy determinant obey that precise
one-dimensional identity. The bridge is a finite boundary condition for
orthogonal polynomials, which selects three Miwa shifts with the required
integer displacements. The Cauchy determinant and Hirota--Miwa identity
are classical ingredients \cite{OhtaHirotaTsujimotoImai93}; the proof below
is self-contained and does not assert that these general principles are new.

\begin{theorem}[Cigler's boundary Hankel recurrence]
\label{thm:cigler-boundary}
For every integer $k\geq0$, every integer $n\geq2k+2$, and all complex
parameters $x,c$,
\begin{align}
 D_{k,n}D_{k,n-2k-2}
 -x^2D_{k,n-1}D_{k,n-2k-1}
 +(x^2-1)D_{k,n-k-1}^{\,2}=0.
 \label{eq:cigler-target}
\end{align}
All determinants in this formula have the same parameters $(x,c)$.
In particular, the unchanged recurrence in Cigler's Conjecture~4.1 holds,
including parameter specializations and indices where determinants vanish.
\end{theorem}

\subsection{The auxiliary determinant problem}

We first derive the precise determinant identity needed for the transfer.
Let $p_i,q_i,w_i$ be parameters, $1\leq i\leq s$, with
$p_i\ne q_j$. Set
\[
 C_{ij}=\frac{(p_i-q_i)w_i}{p_i-q_j},\qquad T=\det(I+C).
\]
For a direction $a$, write $T_a$ for the result of replacing each $w_i$
by $w_i(a-p_i)/(a-q_i)$. Multiple subscripts mean that these replacements
are combined.

\begin{lemma}[Three-direction Cauchy identity]\label{lem:cauchy-three}
For distinct directions $a,b,c$, avoiding all $q_i$,
\begin{equation}\label{eq:cauchy-HM}
 (a-b)T_{ab}T_c+(b-c)T_{bc}T_a+(c-a)T_{ac}T_b=0.
\end{equation}
\end{lemma}
\begin{proof}
Put $v_i=(p_i-q_i)w_i$, $M=I+C$, and
$r_a=(1/(a-q_i))_{i=1}^s$. We work first where $M$ is invertible;
the resulting rational identity then extends wherever its expressions
are defined. Write $P=\operatorname{diag}(p_i)$ and set
\[
 F_a=1-r_a^{\mathsf T}M^{-1}v,\qquad
 G_a=a-r_a^{\mathsf T}M^{-1}Pv.
\]
Conjugating the matrix defining $T_a$ by
$\operatorname{diag}(a-q_i)$ changes it to $M-vr_a^{\mathsf T}$, since
\[
 \frac{a-p_i}{(p_i-q_j)(a-q_j)}
       =\frac1{p_i-q_j}-\frac1{a-q_j}.
\]
The matrix determinant lemma therefore gives $T_a/T=F_a$.
For the two-direction shift the same conjugation, now by
$\operatorname{diag}((a-q_i)(b-q_i))$, gives
\[
 M+Pv\,u_{ab}^{\mathsf T}+v\,v_{ab}^{\mathsf T},\qquad
 u_{ab}=\frac{r_b-r_a}{a-b},\qquad
 v_{ab}=\frac{b r_a-a r_b}{a-b}.
\]
Indeed, subtracting $1/(p_i-q_j)$ from the transformed Cauchy entry
leaves $(p_i+q_j-a-b)/((a-q_j)(b-q_j))$ times $v_i$.
Applying the rank-two determinant lemma and expanding its $2\times2$
determinant yields
\[
 (a-b)\frac{T_{ab}}T=G_aF_b-G_bF_a.
\]
After division by $T^2$, the left side of \eqref{eq:cauchy-HM} is
\[
 (G_aF_b-G_bF_a)F_c+(G_bF_c-G_cF_b)F_a
                         +(G_cF_a-G_aF_c)F_b=0.
\]
This proves the identity without any assumption that a shifted determinant
is nonzero.
\end{proof}

\begin{lemma}[Boundary constraint selecting integer shifts]
\label{lem:cigler-spectral-closure}
Fix $K\geq2$ and $x\ne0,\pm1$. Suppose the nonzero parameters
$\lambda_1,\ldots,\lambda_s$ satisfy
\begin{equation}\label{eq:cigler-dispersion}
 \lambda_i^K-\lambda_i^{-K}
       =x(\lambda_i^{K-1}-\lambda_i^{-(K-1)}).
\end{equation}
Assume the following determinants and direction shifts are defined.
For arbitrary amplitudes $w_i$, let
\[
 \tau_m=\det\left(\delta_{ij}
       +\frac{(\lambda_i-\lambda_i^{-1})w_i\lambda_i^{2m}}
                    {\lambda_i-\lambda_j^{-1}}\right)_{i,j=1}^s.
\]
Then, for every integer $m$,
\begin{equation}\label{eq:cigler-tau-closure}
 \tau_{m+K}\tau_{m-K}
       -x^2\tau_{m+K-1}\tau_{m-K+1}
       +(x^2-1)\tau_m^2=0.
\end{equation}
\end{lemma}
\begin{proof}
In Lemma~\ref{lem:cauchy-three}, use
$p_i=\lambda_i$, $q_i=\lambda_i^{-1}$ and the three directions
$a=0$, $b=x$, $c=x^{-1}$. The corresponding amplitude ratios are
\[
 \frac{0-\lambda_i}{0-\lambda_i^{-1}}=\lambda_i^2,
 \qquad
 \frac{x-\lambda_i}{x-\lambda_i^{-1}}=\lambda_i^{-2(K-1)},
 \qquad
 \frac{x^{-1}-\lambda_i}{x^{-1}-\lambda_i^{-1}}=\lambda_i^{2K}.
\]
The last two identities follow by rearranging
\eqref{eq:cigler-dispersion} as
\[
 \lambda_i^{2K-1}(\lambda_i-x)=1-x\lambda_i.
\]
Thus the three directions increase the determinant index by
$1$, $-(K-1)$, and $K$, respectively. The three-direction identity becomes
\[
 -x\,\tau_{r-K+2}\tau_{r+K}
 +(x-x^{-1})\tau_{r+1}^2
 +x^{-1}\tau_{r+K+1}\tau_{r-K+1}=0.
\]
Multiplication by $x$ and the substitution $m=r+1$ give
\eqref{eq:cigler-tau-closure}.
\end{proof}

The auxiliary statement allows arbitrary amplitudes, not merely the
special coefficients coming from Cigler's moments. Its role is to turn
an exact spectral boundary relation into an exact nonlinear determinant
identity; numerical moment matching is not used.

\subsection{From the path moments to the constrained determinant}

\begin{proof}[Proof of Theorem~\ref{thm:cigler-boundary}]
First remove $c$. If $J_c$ is the infinite symmetric tridiagonal matrix
with diagonal $(c+x,c,c,\ldots)$ and both off-diagonals equal to one,
then $a_{m,k}(x,c)=(J_c^m)_{0k}$. Only a finite portion of this matrix
is involved in each entry. Since $J_c=cI+J_0$,
\[
 a_{m,k}(x,c)=\sum_{j=0}^{m}\binom mj c^{m-j}a_{j,k}(x,0).
\]
Vandermonde's binomial convolution shows that the size-$n$ Hankel matrix
at $c$ equals $B_n$ times the matrix at $0$ times $B_n^{\mathsf T}$,
where $(B_n)_{ij}=\binom ij c^{i-j}$ for $i\geq j$ and zero otherwise.
The matrix $B_n$ is unitriangular, so
\begin{equation}\label{eq:cigler-c-independence}
 D_{k,n}(x,c)=D_{k,n}(x,0).
\end{equation}

Let $\mathcal L(f)=(f(J_0))_{00}$ and define monic polynomials
\[
 p_0(t)=1,\qquad p_1(t)=t-x,\qquad
 p_{j+1}(t)=t p_j(t)-p_{j-1}(t)\quad(j\geq1).
\]
The recurrence gives $p_j(J_0)e_0=e_j$, and therefore
\[
 \mathcal L(p_i p_j)=\delta_{ij},\qquad
 \mathcal L(t^m p_k(t))=a_{m,k}(x,0).
\]
For $k=0$, the unit monic norms give $D_{0,n}=1$ for every $n$.
Then \eqref{eq:cigler-target} is $1-x^2+(x^2-1)=0$; this familiar
case is included only for completeness. Henceforth assume $k\geq1$.
Thus $D_{k,n}$ is the Hankel determinant of the polynomial modification
$p_k\mathcal L$ of an orthonormal functional with all monic norms equal
to one. The classical Christoffel determinant identity, in the form of
\cite[Theorem~1]{Krattenthaler23}, gives
\begin{equation}\label{eq:cigler-Christoffel}
 D_{k,n}(x,0)
    =(-1)^{nk}\frac{\det(p_{n+j}(\rho_i))_{1\leq i\leq k,\,0\leq j<k}}
                        {\Delta(\rho)},
 \qquad \Delta(\rho)=\prod_{i<j}(\rho_j-\rho_i),
\end{equation}
where $\rho_i$ are the roots of $p_k$, initially assumed distinct.
Here is a formal-algebraic derivation, so no analytic measure is
assumed. For $y=(y_1,\ldots,y_n)$, expansion of the two determinants
in the Vandermonde square gives the formal Andr\'eief identity
\[
 D_{k,n}=\frac1{n!}\mathcal L^{\otimes n}
       \left[\Delta(y)^2\prod_{j=1}^n p_k(y_j)\right].
\]
The combined Vandermonde satisfies
\[
 \Delta(y,\rho)=(-1)^{nk}\Delta(y)\Delta(\rho)
                         \prod_{j=1}^n p_k(y_j).
\]
Replace its monomial columns by $p_0,p_1,\ldots,p_{n+k-1}$, and
replace those of $\Delta(y)$ by $p_0,\ldots,p_{n-1}$; monicity leaves
both determinants unchanged. Expand the combined determinant along its
first $n$ rows. Orthogonality makes every term vanish unless the selected
columns are exactly $p_0,\ldots,p_{n-1}$. That surviving term has positive
Laplace sign, and applying $\mathcal L^{\otimes n}$ contributes $n!$.
The remaining minor is $(p_{n+j}(\rho_i))$. Substitution proves
\eqref{eq:cigler-Christoffel}, with its stated sign and denominator.
For $n=0$ the identity reduces directly to monicity and both sides equal one.

Work now over generic $x$, choose
$\rho_i=\lambda_i+\lambda_i^{-1}$, and put
\[
 A_i=\frac{\lambda_i-x}{\lambda_i-\lambda_i^{-1}},\qquad
 B_i=\frac{x-\lambda_i^{-1}}{\lambda_i-\lambda_i^{-1}}.
\]
The explicit solution of the three-term recurrence is
\[
 p_j(\rho_i)=A_i\lambda_i^j+B_i\lambda_i^{-j}.
\]
The boundary condition $p_k(\rho_i)=0$ says
$B_i=-A_i\lambda_i^{2k}$ and also gives
\eqref{eq:cigler-dispersion} with $K=k+1$.

Define
\[
 R_i=\prod_{j\ne i}
       \frac{\lambda_i-\lambda_j^{-1}}
            {\lambda_i^{-1}-\lambda_j^{-1}},
 \qquad w_i=\frac{A_i}{B_i}R_i.
\]
Expanding the rows of the determinant in \eqref{eq:cigler-Christoffel}
into their two exponential terms gives
\begin{align}
 \det(p_{n+j}(\rho_i))
 &=\left(\prod_i B_i\right)\Delta(\lambda^{-1})
                    \left(\prod_i\lambda_i\right)^{-n}\tau_n,
 \label{eq:cigler-soliton-gauge}\\
 \tau_n
 &=\sum_{S\subseteq\{1,\ldots,k\}}
       \prod_{i\in S}(w_i\lambda_i^{2n})
       \prod_{\substack{i<j\\i,j\in S}}
          \left(\frac{\lambda_i-\lambda_j}{1-\lambda_i\lambda_j}\right)^2.
 \notag
\end{align}
Here $\Delta(\lambda^{-1})=\prod_{i<j}(\lambda_j^{-1}-\lambda_i^{-1})$.
To verify the expansion, first factor the term in which every row uses
$B_i\lambda_i^{-n-j}$. Switching row $i$ alone contributes
$(A_i/B_i)\lambda_i^{2n}R_i$. Switching both $i$ and $j$ corrects
these individual contributions by the displayed squared factor.
All Vandermonde factors are products over pairs, so these singleton and
pair corrections account for every subset.
The Cauchy determinant formula shows that this subset sum is exactly
\[
 \tau_n=\det\left(\delta_{ij}
       +\frac{(\lambda_i-\lambda_i^{-1})w_i\lambda_i^{2n}}
                    {\lambda_i-\lambda_j^{-1}}\right)_{i,j=1}^k:
\]
the principal minor indexed by $S$ is its corresponding summand.
Consequently Lemma~\ref{lem:cigler-spectral-closure} applies.

Equations \eqref{eq:cigler-Christoffel} and
\eqref{eq:cigler-soliton-gauge} give
\[
 D_{k,n}(x,0)=Cz^n\tau_n,\qquad
 C=\frac{(\prod_i B_i)\Delta(\lambda^{-1})}{\Delta(\rho)},
 \qquad z=\frac{(-1)^k}{\prod_i\lambda_i}.
\]
Every product in the centered identity
\eqref{eq:cigler-tau-closure} acquires the same factor $C^2z^{2m}$.
Hence that identity holds with $\tau_j$ replaced by $D_{k,j}$.
Set $m=n-k-1$; the three products then have the exact indices in
\eqref{eq:cigler-target}.

Finally, the calculation is valid on a nonempty generic set of $x$.
Indeed at $x=0$ the roots of $p_k=U_k(t/2)$ are distinct and lie in
$(-2,2)$, so none equals $2$ or $-2$; nearby generic nonzero $x$ avoids
all the finitely many denominator conditions. For fixed $k,n$, the left
side of \eqref{eq:cigler-target} is a polynomial in $x$ and vanishes on
this generic set. It is therefore the zero polynomial. This proves the
identity also at $x=0,\pm1$, at root collisions, and at every zero of any
Hankel determinant. Equation~\eqref{eq:cigler-c-independence} restores
arbitrary $c$. No division by an original Hankel determinant is used in
the final extension.
\end{proof}

\paragraph{Scope and attribution.}
The theorem proves the entire stated recurrence, not merely the examples
$k=1,2,3,4$, and not a determinant sequence with changed moments.
The method passes through a new auxiliary spectral-closure question,
but Christoffel modification, Cauchy minors, and the Hirota--Miwa equation
are prior mathematics. The specific boundary-to-shift identification is
the proposed contribution. A bounded search found no earlier proof of
this exact Cigler recurrence; that does not certify universal historical
priority or a historically unprecedented foundational principle.

\section{A complete two-level tree classification and an excluded prior target}
\label{sec:two-level-trees}

The original target considered here was Open Problem~7.1 of
Coutinho--Guo--Schmeits \cite{CGS25}: whether infinitely many non-star
finite trees admit Laplacian peak state transfer. A public construction
of such trees was posted on 9 October 2026 and merged on 10 October in
\cite{PriorTree26}. Therefore this existence statement is not counted as
a newly solved open problem here. We nevertheless obtain a complete
classification of every transferring vertex pair in the following
natural two-parameter family. This is an auxiliary result, not a
replacement for the original success criterion.

For integers $a\geq1$ and $b\geq2$, let $T(a,b)$ have a root $o$, children
$c_1,\ldots,c_b$, and $a$ private pendant leaves $\ell_{i,1},\ldots,
\ell_{i,a}$ attached to $c_i$. It is a non-star tree on
$N=1+b+ab$ vertices. For a real symmetric matrix $M=\sum_\lambda
\lambda E_\lambda$ and distinct vertices $x,y$, nonzero peak transfer
means
\[
 |(e^{itM})_{xy}|=\sum_\lambda |(E_\lambda)_{xy}|>0.
\]
Equivalently, all nonzero summands $e^{it\lambda}(E_\lambda)_{xy}$
have the same complex phase. For a connected graph Laplacian,
$(E_0)_{xy}=1/N>0$, so that common phase must be $1$.

\begin{theorem}[Classification of Laplacian peaks in two-level trees]
\label{thm:two-level-classification}
Let $a\geq1$, $b\geq2$. Every distinct pair admitting Laplacian peak
transfer in $T(a,b)$ contains the root $o$. Such a pair exists if and
only if there are integers $u\geq1$, $v\geq2$ of opposite parity such
that
\[
 a=uv,\qquad b=uv-u+v.
\]
These parameters, when they exist, are unique. Set
\[
 r=u(v-1)+1,\qquad s=v(u+1)+1,\qquad g=\gcd(r,s).
\]
Then the entire list of transferring pairs is as follows:
\begin{enumerate}
\item If $u$ is odd and $v$ is even, the root has peak transfer to
      every child, and to no leaf.
\item If $u$ is even and $v$ is odd, the root has peak transfer to
      every leaf, and to no child.
\end{enumerate}
In either case the positive peak times are exactly the odd multiples
of $\pi/g$. The peak amplitudes for the two cases are respectively
\[
 B_L(o,c_i)=\frac{2(s-1)}{s(s-r)},\qquad
 B_L(o,\ell_{i,j})=\frac{2}{r(s-r)}.
\]
If the indicated parametrization with opposite parity does not exist,
there is no Laplacian peak transfer between any distinct vertices.
Equivalently, existence of a transferring pair is characterized by
$a$ even, $b$ odd, and $(a-b)^2+4a$ being a perfect square.
\end{theorem}

\begin{proof}
The normalized indicators of the root, all children, and all leaves
span an invariant subspace for the Laplacian. Its matrix is
\[
 Q=\begin{pmatrix}
 b&-\sqrt b&0\\
 -\sqrt b&a+1&-\sqrt a\\
 0&-\sqrt a&1
 \end{pmatrix},\qquad
 p(z):=\det(zI-Q)=z\bigl(z^2-(a+b+2)z+N\bigr).
\]
Let $0<r<s$ be the two positive roots. Since
\[
 (b-r)(b-s)=1-b<0,
\]
we have $r<b<s$. Also the quadratic has positive value $a(b-1)$
at $1$ and negative value at $b>1$, so its smaller root lies in $(1,b)$.

Since the root vector lies in this invariant subspace, it has zero
projection onto its orthogonal complement. Dividing the quotient
resolvent entries by the appropriate cell normalization factors gives
for each individual child and leaf
\[
 ((zI-L)^{-1})_{o,c_i}=-\frac{z-1}{p(z)},\qquad
 ((zI-L)^{-1})_{o,\ell_{i,j}}=\frac1{p(z)}.
\]
All three residues are nonzero. At the ordered eigenvalues $0,r,s$,
the child coefficients have signs $(+,+,-)$, and the leaf coefficients
have signs $(+,-,+)$. Thus the mutual support is exactly $\{0,r,s\}$.

If a peak occurs, the phases at $r$ and $s$ must both be real, so
$tr,ts\in\pi\mathbb Z$. Hence $r/s$ is rational. As $r+s=a+b+2$
is integral, both roots are rational; they are roots of a monic
integer polynomial, so they are integers. Conversely, for integral
$r,s$, write $r=gr_0$, $s=gs_0$ with $\gcd(r_0,s_0)=1$.
The integer combinations of $r_0,s_0$ show that $tg\in\pi\mathbb Z$
is necessary. The child sign pattern can occur exactly when $r_0$
is even and $s_0$ is odd; the leaf sign pattern requires $r_0$ odd and
$s_0$ even. In each admissible case the allowed positive multiples of
$\pi/g$ are exactly the odd ones.

It remains to express integrality and parity in the original graph
parameters. The discriminant is
\[
 D=(a-b)^2+4a.
\]
If it is a square $h^2$, then $h$ and $a-b$ have the same parity and
$h>|a-b|$. Therefore
\[
 u=\frac{h+a-b}{2},\qquad v=\frac{h-a+b}{2}
\]
are positive integers with $uv=a$ and $b=uv-u+v$. The condition
$b\geq2$ excludes $v=1$. Conversely these formulas give
$D=(u+v)^2$ and the roots stated in the theorem. They also establish
uniqueness. At least one of $r,s$ is odd: checking the four parities
of $u,v$ gives
\[
\begin{array}{c|c|c}
(u\bmod2,v\bmod2)&r\bmod2&s\bmod2\\\hline
(1,0)&0&1\\
(0,1)&1&0\\
(1,1)&1&1\\
(0,0)&1&1.
\end{array}
\]
Thus $g$ is odd and the required root parities are exactly the two
opposite-parity cases in the theorem. Since the sum of all cross
projector coefficients is zero, the peak amplitude is twice the
absolute value of the unique negative coefficient. The residues give
$-(s-1)/(s(s-r))$ for a child and $-1/(r(s-r))$ for a leaf, proving
the stated amplitudes.

Finally consider two distinct nonroot vertices. Besides the radial
subspace above, the Laplacian has a branch-contrast subspace. Choose
$t=(t_1,\ldots,t_b)$ with $\sum_i t_i=0$. Vectors with child values
$xt_i$ and all leaf values $yt_i/\sqrt a$ in branch $i$ reduce the
Laplacian to
\[
 R=\begin{pmatrix}a+1&-\sqrt a\\-\sqrt a&1\end{pmatrix}.
\]
Its eigenvalues $\rho_-<\rho_+$ are the two roots of
$z^2-(a+2)z+1$. They are irrational, because
\[
 (a+1)^2<a^2+4a<(a+2)^2\qquad(a\geq1).
\]
Both coordinates of either normalized eigenvector of $R$ are nonzero.
The remaining leaf-contrast subspaces consist of vectors summing to zero
within a private leaf set and have eigenvalue $1$. The dimensions
$3+2(b-1)+b(a-1)=N$ account for the full vertex space.
For nonroot vertices in branches $i,j$, the cross entry of the
$\rho_\pm$ projector is the product of these nonzero coordinates
(with a factor $1/\sqrt a$ per leaf), multiplied by
$\delta_{ij}-1/b$. This factor never vanishes for $b\geq2$.
There is no eigenvalue collision with the radial subspace: subtracting
the two quadratics would force a common root to equal $a+1$, whereas
$R$'s characteristic polynomial at $a+1$ equals $-a$.
Nor is either root equal to the leaf-difference eigenvalue $1$.
Consequently both $\rho_-,\rho_+$ lie in the mutual support of every
nonroot pair. The zero eigenvalue also lies in that support.
Peak alignment would imply $\rho_+/\rho_-\in\mathbb Q$, which,
combined with $\rho_++\rho_-=a+2$, would make both roots rational.
This contradiction excludes every nonroot pair and completes the
classification.
\end{proof}

Taking $u=1$, $v=2k$ gives the particularly small infinite family
\[
 a=2k,\qquad b=4k-1,\qquad N=2k(4k+1),\qquad k\geq1,
\]
with root-to-child peak transfer at time $\pi$. Its first orders are
$10,36,78,136,\ldots$. Taking $u=2$, $v=3$ gives a $50$-vertex tree
with root-to-leaf peaks at time $\pi/5$. The public prior family for
the original existence question is instead $a=t(t+1)$,
$b=a+1$, for positive odd $t$, with roots $t^2+1$ and $(t+1)^2+1$.
That prior construction already settles the unchanged existence
question. The classification above supplies more detail within a
restricted family, but does not constitute a first resolution of
Open Problem~7.1 or a classification of all trees.

The proof uses equitable subspaces, spectral projectors, and elementary
integer factorization. These are classical tools. A historically new
foundational method is not asserted.

\appendix

\section{Least dimension: an unsuccessful projection route}
\label{sec:ehrhart-projection-audit}

Clarke's Question~6.8 \cite{Clarke26} asks for the least dimension of an
order-two counterexample to polynomiality implying effectiveness. We do
\emph{not} resolve it here. Combining the source's six-dimensional example
with the preceding T83 internal proof \cite{Keio81} gives only
$4\leq d_{\min}\leq6$. The obstruction below excludes a particular route
to dimension five; it is not counted as a new named-problem solution.

Let $P_0\subset\mathbb R^6$ be the convex hull of
\[
\begin{array}{llll}
(1,0,0,0,0,0),&(0,0,1,0,0,0),&(0,1,0,1,1,0),&(0,0,1,1,1,0),\\
(0,0,0,1,0,0),&(0,0,0,0,0,1),&(1,1,0,0,1,0),&(1,1,0,0,0,1).
\end{array}
\]
This is Clarke's Example~5.6, with involution $\sigma(u,v)=(v,u)$.
For $a,b\in\mathbb Z$, define the surjective lattice quotient
\[
 T_{a,b}=
 \begin{pmatrix}
 1&0&0&0&0&a\\0&1&0&0&0&b\\0&0&1&0&0&1\\
 0&0&0&1&0&-a\\0&0&0&0&1&-b
 \end{pmatrix},\qquad Q_{a,b}=T_{a,b}(P_0).
\]
Its kernel is generated by the primitive anti-invariant vector
$(a,b,1,-a,-b,-1)$. The induced involution is
\[
 (y_1,y_2,y_3,y_4,y_5)\longmapsto
 (y_4+ay_3,y_5+by_3,y_3,y_1-ay_3,y_2-by_3).
\]
Consequently $Q_{a,b}$ is five-dimensional, and the homogeneous
plus/minus eigenspace dimensions are $(p,q)=(4,2)$.

\paragraph{Projection obstruction.}
For every $a,b\in\mathbb Z$, the equivariant numerator of $Q_{a,b}$
is not polynomial.

\begin{proof}
The fixed lattice has coordinates $(y_4,y_5,y_3)$. Averaging the vertices
shows that twice the fixed polytope $R_{a,b}=Q_{a,b}^{\sigma}$ is
\[
 2R_{a,b}=\operatorname{conv}\{(-a,-b,2),(1,0,0),
                                  (1-a,1-b,2),(1,2,0)\}.
\]
Changing $(a,b)$ to $(a+2r,b+2s)$ gives the lattice shear
$(x,y,z)\mapsto(x-rz,y-sz,z)$. Moreover the affine lattice involution
$(x,y,z)\mapsto(1-x-z,1-y-z,z)$ sends $R_{0,0}$ to $R_{1,1}$ and
$R_{0,1}$ to $R_{1,0}$. Thus there are only two counting cases.

Here is an explicit finite certificate, independent of any hull software.
Put
\[
 A_{k,z}=\max\{0,\lfloor k/2\rfloor-\lceil(k-z)/2\rceil+1\}.
\]
The four supporting inequalities of $R_{0,0}$ give, in $kR_{0,0}$,
\[
 0\leq z\leq k,\quad (k-z)/2\leq x\leq k/2,\quad
 x-(k-z)/2\leq y\leq x+(k-z)/2.
\]
For $kR_{0,1}$, the $x,z$ bounds are identical and the $y$ bounds become
$x-k/2\leq y\leq x-z+k/2$. Therefore their lattice counts are exactly
\begin{align*}
 N_{0,0}(k)&=\sum_{z=0}^{k}A_{k,z}\bigl(2\lfloor(k-z)/2\rfloor+1\bigr),\\
 N_{0,1}(k)&=\sum_{z=0}^{k}A_{k,z}\bigl(2\lfloor k/2\rfloor-z+1\bigr).
\end{align*}
For $k=0,\ldots,7$ these sums give
\[
\begin{array}{c|rrrrrrrr}
 k&0&1&2&3&4&5&6&7\\\hline
 N_{0,0}(k)&1&1&6&6&19&19&44&44\\
 N_{0,1}(k)&1&0&7&3&22&13&50&34
\end{array}
\]
Each $R_{a,b}$ has dimension three and denominator at most two, so its
even and odd Ehrhart constituents have degree at most three. Four values
of each parity determine them, giving
\[
\begin{array}{c|cc}
 &a\equiv b\pmod2&a\not\equiv b\pmod2\\\hline
 N_{\mathrm{even}}(k)&
 \dfrac{(k+2)(k^2+4k+6)}{12}&
 \dfrac{(k+2)(k+4)(2k+3)}{24}\\[5pt]
 N_{\mathrm{odd}}(k)&
 \dfrac{(k+1)(k^2+2k+3)}{12}&
 \dfrac{(k-1)(k+1)(2k+3)}{24}\\[5pt]
 N_{\mathrm{even}}-N_{\mathrm{odd}}&
 \dfrac{k^2+3k+3}{4}&\dfrac{(2k+3)^2}{8}.
\end{array}
\]
In both cases the difference has degree two, hence the fixed Ehrhart
series has a pole of order three at $t=-1$. The character formula
\[
 \varphi_{Q_{a,b}}(t)(\sigma)
  =(1-t)^4(1+t)^2\sum_{k\geq0}N_{a,b}(k)t^k
\]
cancels only two orders. This proves the claimed nonpolynomiality for all
integer parameters, not just those in a bounded search.
\end{proof}

\paragraph{A rejected four-dimensional candidate.}
Under $(x_1,x_2,y_1,y_2)\mapsto(y_1,y_2,x_1,x_2)$, consider
\[
 P=\operatorname{conv}\{(1,0,0,0),(0,0,1,0),(1,1,0,1),
              (0,1,1,1),(2,1,1,0),(1,0,2,1)\}.
\]
Its fixed triangle is
$\operatorname{conv}\{(1/2,0),(1/2,1),(3/2,1/2)\}$, with constituents
$N_{\mathrm{even}}(k)=k^2/2+k+1$ and
$N_{\mathrm{odd}}(k)=(k^2-1)/2$.
Exact counts $1,6,23,68,165$ at dilations $0,\ldots,4$ give
$h^*_P=1+t+3t^2+3t^3$, and consequently
\[
 \varphi_P=(1+3t^2+2t^3)\chi_0+(t+t^3)\chi_1.
\]
It is effective and therefore fails as a counterexample. The accompanying
exact-arithmetic certificate supplies its supporting facets and verifies
all the displayed counts. Neither this rejection nor the projection
obstruction supplies the missing four- and five-dimensional classification.

\section{A sparse-PT construction that cannot solve the odd-period problem}

Faust--Friedman--O'Malley--Ramos--Sharma, Conjecture~2 in
\emph{Algebraic Properties of the Ideal of Spectral Invariants for the
Discrete Laplacian} \cite{Faust26}, asks for a nonzero reversal-conjugate
$q$-periodic potential Floquet isospectral to zero for every $q>3$.
Their source proves every even period. The uniform odd-period assertion
is not proved here. The following proposition rules out one proposed
uniform construction; it is not a solution of that conjecture.

\begin{proposition}[Five-site reflected obstruction]
Let $r\geq1$ and $q=2r+5$. For real $a,b,c,d$, put
\[
 V=(a+ib,c+id,\underbrace{0,\ldots,0}_{r},-2(a+c),
       \underbrace{0,\ldots,0}_{r},c-id,a-ib).
\]
If the nearest-neighbor discrete Schr\"odinger operator with potential
$V$ and off-diagonal weights one is Floquet isospectral to zero, then
$a=b=c=d=0$.
\end{proposition}
\begin{proof}
Write
\[
 M(z)=\begin{pmatrix}x-z&-1\\1&0\end{pmatrix},\quad
 A=M(0),\quad C=M(a+ib),\quad B=M(c+id),\quad h=-2(a+c),
\]
and put $D=BC\overline C\,\overline B$, where conjugation fixes $x$.
By cyclicity of trace, the transfer discriminant of $V$ is
$\operatorname{tr}(A^rM(h)A^rD)$; that of zero is
$\operatorname{tr}(A^{2r+5})$. Set
$u=U_r(x/2)$ and $v=U_{r-1}(x/2)$ for the Chebyshev polynomials of
second kind. Then
\[
 A^r=\begin{pmatrix}u&-v\\v&u-xv\end{pmatrix}.
\]
The discriminant difference has the form
\[
 F(x)u^2+G(x)uv+H(x)v^2,
 \qquad H(x)=(-3a^2+b^2)x+2a(a^2+b^2).
\]
The formula for $H$ follows by substituting $u=0,v=1$ in the matrix
expression; it is a polynomial identity independent of $r$.
Consecutive Chebyshev polynomials are coprime. Consequently a vanishing
discriminant difference implies $u\mid H$.
If $r=1$, then $u=x$, so the constant coefficient of $H$ vanishes.
If $r\geq2$, then $\deg u\geq2>\deg H$, so $H=0$.
In either case $a(a^2+b^2)=0$, and reality gives $a=0$.

The coefficients of $x^{q-2}$ and $x^{q-3}$ in the discriminant
difference are respectively
\[
 -3a^2-4ac+b^2-3c^2+d^2,
 \qquad
 2(a^3+4a^2c+ab^2+4ac^2+c^3+cd^2).
\]
These follow directly from the first three trace identities, with
$\sum_j V_j=0$. Substituting $a=0$ into the second coefficient gives
$c(c^2+d^2)=0$, hence $c=0$. The first coefficient is then $b^2+d^2$,
so $b=d=0$.
\end{proof}

The central value was imposed by the trace, so the proposition exhausts
this five-site reflected pattern. Period five corresponds to $r=0$;
then $u=1$ gives no restriction, consistently with the nonzero solutions
already certified in the source. Exact symbolic calculations verify the
transfer identity and coefficient formulas for $q=7,9,11,13,15,17$;
the divisibility argument proves the assertion for every $r\geq1$.

This removes a candidate construction, but does not imply nonexistence
of less sparse PT potentials. The original all-odd-period conjecture and
the unrestricted period-lattice classification remain unresolved in this
attempt. Transfer matrices, Chebyshev identities and polynomial
coprimality are classical ingredients; no historically new foundational
principle is claimed.

\section{Verification, attribution, and the success threshold}

The accompanying archive supplies seven executable Python verification
programs and their recorded JSON outputs. Each uses exact integer,
rational, or symbolic arithmetic. Finite tests check implementation and
formula consistency; the preceding proofs establish the infinite
quantifiers. The distinction is especially important for the Hankel
recurrence, whose early zero blocks can make short tests vacuous.

\begin{center}\small
\begin{tabular}{@{}>{\raggedright\arraybackslash}p{.25\textwidth}>{\raggedright\arraybackslash}p{.67\textwidth}@{}}
\toprule
Component & Executed checks \\
\midrule
Stern families $0,2,3$ & Eight exact Bernstein identities; 177 independent Stern-recursion comparisons, covering all three families for $2\leq n\leq60$. \\
Fourth Stern family & Sixty exact factorization checks; twelve exact Sturm and multiplicity checks. The proof separately covers every $n$. \\
Cigler recurrence & 108 parameter runs comparing direct Hankel determinants with a fixed-size reduction; 24 long-range recurrence checks for $1\leq k\leq12$, $x=2,3$, reaching $n=512$. \\
Cauchy identity & Exact rational Cauchy-minor expansions and nontrivial three-direction identities for one through six spectral pairs. \\
Tree classification & 870 parameter cases and seven full Laplacian matrices of trees with at most 50 vertices. Individual vertex projectors, not only quotient states, are checked. \\
PT obstruction & Exact symbolic transfer discriminants and coefficient identities for periods $7,9,11,13,15,17$. The Chebyshev divisibility proof covers all odd periods at least seven in the stated pattern. \\
Ehrhart attempts & Exact supporting-facet and lattice-point calculations for the four-dimensional candidate and the four parity classes of the projected fixed polytopes. \\
\bottomrule
\end{tabular}
\end{center}

Independent internal reviews rederived the Stern identities and all
positivity certificates, and audited the fourth-family root counts.
For the Hankel theorem, they independently checked the Cauchy rank-two
identity, reciprocal-root shifts, Vandermonde normalization, and the
extension to exceptional parameters. An alternative Casoratian--Pl\"ucker
derivation was also reconstructed during the audit. These reviews are
not external peer review or formal proof-assistant certification.

The original named targets that receive complete, unambiguous arguments
in this manuscript are Ulas Conjecture~4.3(1) and Cigler Conjecture~4.1.
The fourth Stern family is fully classified, but the multiplicity
ambiguity of its source statement prevents treating it as a separate
unqualified historical closure. The already preceded tree existence
question, the two unresolved main questions, and the initial Motzkin
shifted-period attempt are excluded from this count.

The priority search inspected the original papers, author-hosted
material, follow-up literature, and targeted public proof records through
26 October 2026. No earlier identical complete proof of the two principal
applications was located. This is the outcome of a bounded search, not
a theorem about all prior literature. The general methods have classical
antecedents, explicitly credited in the body and bibliography. Thus the
full demand for five historically unprecedented-method solutions is not
reported as achieved.

The standalone TEX embeds its bibliography and uses no external figures
or bibliography database. It compiles with two passes of
\texttt{pdflatex}. The archive includes dependency information, run
instructions, and SHA-256 checksums. Recorded calculation outputs are
provided to make the finite checks inspectable and reproducible.


\begin{thebibliography}{99}
\bibitem{Ulas20}
M. Ulas, \emph{Strong arithmetic property of certain Stern polynomials},
Publicationes Mathematicae Debrecen \textbf{96}(3--4) (2020), 401--422.
\href{https://doi.org/10.5486/PMD.2020.8683}{doi:10.5486/PMD.2020.8683}.
Inspected preprint: \href{https://arxiv.org/abs/1909.10844v1}{arXiv:1909.10844v1} (2019).

\bibitem{CiglerRemarks}
J. Cigler, \emph{Some remarks and conjectures about Hankel determinants},
version 5 (2010), Conjecture~4.1, equation~(4.5).
\href{https://homepage.univie.ac.at/johann.cigler/preprints/conjectures-hankel.pdf}{Author's preprint};
\href{https://homepage.univie.ac.at/johann.cigler/prepr.html}{author's dated index}.

\bibitem{CiglerMO23}
J. Cigler, \emph{A conjecture about Hankel determinants of path generating functions},
MathOverflow question 456683, posted 18 October 2023, edited 22 October 2023.
\href{https://mathoverflow.net/questions/456683}{Original question}.

\bibitem{Krattenthaler23}
C. Krattenthaler, \emph{Hankel determinants of linear combinations of
moments of orthogonal polynomials, II}, Ramanujan Journal \textbf{61}
(2023), 597--627.
\href{https://arxiv.org/abs/2101.04225v5}{arXiv:2101.04225v5} (2021).

\bibitem{OhtaHirotaTsujimotoImai93}
Y. Ohta, R. Hirota, S. Tsujimoto, and T. Imai,
\emph{Casorati and discrete Gram type determinant representations of
solutions to the discrete KP hierarchy}, Journal of the Physical Society
of Japan \textbf{62}(6) (1993), 1872--1886.
\href{https://doi.org/10.1143/JPSJ.62.1872}{doi:10.1143/JPSJ.62.1872}.

\bibitem{CGS25}
G. Coutinho, K. Guo, and V. Schmeits,
\emph{Peak state transfer in continuous quantum walks},
\href{https://arxiv.org/abs/2505.11986v4}{arXiv:2505.11986v4},
9 September 2026 (first preprint 2025), Open Problem~7.1.

\bibitem{PriorTree26}
The Omega Institute, \emph{Prove infinitely many non-star trees admit
Laplacian peak state transfer (\#14801)}, public pull request 14855,
opened 9 October 2026, merged 10 October 2026.
\href{https://github.com/the-omega-institute/trureturing/pull/14855}{Public construction and proof record}.
The existence of this prior construction is used for attribution;
the tree classification here is proved independently.

\bibitem{Clarke26}
O. Clarke, \emph{The equivariant Ehrhart theory of polytopes with
order-two symmetries: Part II},
\href{https://arxiv.org/abs/2610.17988v1}{arXiv:2610.17988v1},
13 October 2026.

\bibitem{Faust26}
M. Faust, L. Friedman, G. O'Malley, R. Ramos, and A. Sharma,
\emph{Algebraic properties of the ideal of spectral invariants for the
discrete Laplacian},
\href{https://arxiv.org/abs/2602.08304v1}{arXiv:2602.08304v1} (2026).

\bibitem{Keio81}
KEIO, \emph{Derivative obstructions, minimal Ehrhart counterexamples,
and peak-transfer bounds: T81--T85}, unpublished research manuscript,
24 October 2026. The T83 argument excludes three-dimensional
order-two ineffective polynomial numerators; it remains an internally
checked research result.

\end{thebibliography}

\end{document}
